\chapter{水为什么如此特别}

\begin{kaichang}
倒一杯冰水，放进三块冰。现在就盯着它看一分钟，你会看到三件怪事。

第一，冰块浮在水面上。请认真想一想这件事有多怪：一种物质的固体，居然比它自己的液体轻。你把一块固态的蜡扔进熔化的蜡油里，它会沉底；把固态的橄榄油放进液态橄榄油里，它也会沉。几乎所有物质都是固体更密、会下沉。水偏偏反着来。

第二，杯子的外壁慢慢挂上了一层小水珠。杯子里明明没有漏，水是哪里来的？

第三，过半小时再看，冰块化了，水位几乎没变——冰块伸出水面的那一小截，融化后似乎“凭空消失”了。

把这三件怪事记在纸上，再补一个你没见过但猜得到的：为什么发烧时大人用温水给你擦身？为什么住在海边的人，夏天没有内陆那么热、冬天没有内陆那么冷？这些问题的答案，全都藏在一个只有三个原子的小分子里。
\end{kaichang}

\section{一张“反常物质”的罪状清单}

在物理学和化学的教室里，水通常被当作“标准物质”：比热拿它当单位，密度拿它当参照，酸碱中性拿它划线。可是科学家眼里的水，恰恰是\textbf{最反常的常见物质}。先把它的“罪状”一条条列出来：

\begin{itemize}[itemsep=2pt]
\item \textbf{它太爱当液体了}。和水分子“体重”差不多的分子——甲烷（\chem{CH_4}）、氨（\chem{NH_3}）、硫化氢（\chem{H_2S}）——在室温下全是气体。按“亲戚”的脾气推算，水的沸点应该在零下一百度附近，地球上根本不该有江河湖海。可水的沸点偏偏是 $100\,^{\circ}\mathrm{C}$。
\item \textbf{它的固体比液体轻}。冰的密度大约是 \ud{0.917}{g/cm^3}，所以冰块浮在水上，高出水面的部分约占总体积的十分之一——“冰山一角”说的就是这个比例。
\item \textbf{它特别“扛热”}。把同样一壶水和一壶油放在同样的火上，油温升得比水快得多。水每升高 $1\,^{\circ}\mathrm{C}$ 要吞下的热量，几乎是常见液体里最多的。
\item \textbf{它挑食}。盐、糖、酒精在水里一搅就没影了；油却怎么搅都要分家，沙拉酱静置一会儿就分成两层。
\end{itemize}

这些“罪状”之间有没有共同的幕后黑手？有。要抓住它，得把镜头推进到单个水分子的样子——接下来你会发现，这四条看似互不相干的怪脾气，其实全是同一个结构原因在不同场景下的投影。

\section{一个歪歪扭扭、一头重一头轻的分子}

第 19 章讲过，氧和氢之间是极性共价键：氧的电负性大，把共享电子拉向自己，结果氧端带一点负电，氢端带一点正电。第 20 章又讲过，水分子不是直线形，而是弯的——两个氢原子张开的夹角约 $104.5^{\circ}$，像米老鼠的两只耳朵。

这两件事合起来，后果很严重：正电中心和负电中心\textbf{不能重合}。如果水分子是直线形（像二氧化碳那样两头对称），两端的部分电荷会互相抵消，整个分子对外不显电性。可它是弯的，于是“米老鼠的下巴”（氧端）偏负，“耳朵尖”（氢端）偏正。整个水分子就像一粒微小的、歪歪扭扭的磁铁——一头是负极，一头是正极。

接下来发生的事，第 22 章已经埋下过引子：一个水分子的“耳朵尖”（偏正的氢），会被邻居的“下巴”（偏负的氧）吸引，形成一种比化学键弱、但比普通分子间作用强得多的连接——\textbf{氢键}。氢键不是分子内部那种共享电子的共价键，它是分子\textbf{之间}的静电吸引，强度大约只有共价键的二十分之一。

关键在于水分子“长了几只手”。氧原子上还有两对没参与成键的电子（孤对电子），它们可以接受来自另外两个水分子的氢；自己的两个氢又能伸向另外两个邻居。算下来，一个水分子最多可以同时牵住\textbf{四个}邻居。四个方向一张开，恰好指向四面体的四个顶点——这不是巧合，正是那对“米老鼠耳朵”张角接近四面体角的原因。

于是液态水不是一碗各顾各的“分子豆粥”，而是一张不断翻新的大网：每个分子伸出氢键拉住四邻，氢键不断断开又立刻接上，平均寿命只有一万亿分之一秒量级，但任何时刻，绝大多数分子都挂在网里。水的所有怪脾气，都是这张网的投影。

\section{用一张网解释四条罪状}

\subsection{罪状一：沸点高得离谱}

让液体沸腾，本质上是给分子足够的速度，让它们挣脱邻居、飞进空气。甲烷分子之间只有微弱的色散力（第 22 章），轻轻一推就散伙，所以 $-161\,^{\circ}\mathrm{C}$ 就沸腾了。而水分子要飞出去，得先把手上的氢键一根根扯断——这需要多得多的能量，于是水要到 $100\,^{\circ}\mathrm{C}$ 才肯大规模“散伙”。没有氢键的水，是一种地球上不可能存在的假想物质；有了氢键，水才在常温下老老实实做液体，也才有了海洋、云、雨和你身体里的那一大半体重。

\subsection{罪状二：冰浮在水上}

液态时，氢键网是乱糟糟、随时重排的，分子之间有不少“见缝插针”的挤法。结冰时情况变了：温度低，分子的热运动压不过氢键的“取向要求”，每个分子都稳稳牵住四个邻居，搭成一座规整的四面体脚手架。脚手架为了迁就氢键的方向，中间留出了很多空洞——就像为了对齐榫头而搭得稀稀疏疏的架子。结果是：\textbf{冰的结构反而比液态水更空旷}，同样多的分子占了更大的体积，密度变小，自然浮在水面上。水结冰时体积约膨胀 9\%，冬天水管冻裂、岩石被冻胀崩解，都是这 9\% 干的。

这件事对地球生命是个天大的好消息：冬天湖面先结冰，冰浮在表面，像给湖水盖了床保温被，下面的水保持液态，鱼和水草得以过冬。假如冰像别的固体那样下沉，湖泊会从底部向上冻实，北方水体的生态完全是另一副模样。

\subsection{罪状三：特别扛热}

给水加热，热量去哪儿了？在普通液体里，热量大多直接变成分子的乱跑乱撞——温度立刻上升。而在水里，输入的能量先要拿去做一件“家务”：拆断氢键、松开那张网。等网拆松了，剩下的能量才轮到升高温度。所以水的“热容”格外大：每克水升高 $1\,^{\circ}\mathrm{C}$ 需要约 \ud{4.2}{J}，是乙醇的两倍左右，是铁这种金属的约九倍。

反过来，水蒸发时，跑得最快的分子要挣脱整张网才能飞走，顺带卷走一大笔能量——每蒸发 1 克水，大约带走 \ud{2400}{J} 的热量。这就是你出汗和温水擦身能退烧的原因：汗水蒸发，等于把身体的热量“打包快递”给了空气。

这套“扛热＋蒸发散热”的组合，也是你自己的体温调节系统。你的身体约六成是水，这台由肌肉、内脏组成的“发动机”每时每刻都在产热，而体液巨大的热容像一个缓冲水池，让核心体温不会随一顿饭、一阵奔跑大起大落；一旦缓冲池快装不下，汗腺就启动蒸发散热，把多余的热量直接送出体外。可以说，哺乳动物能把体温常年稳在 $37\,^{\circ}\mathrm{C}$ 附近，首先是因为体内灌满了这种“扛热”的液体。

\begin{qiaoliang}
拿一杯水算算账。一杯水约 200 克，把它从 $20\,^{\circ}\mathrm{C}$ 烧到 $100\,^{\circ}\mathrm{C}$，升温 $80\,^{\circ}\mathrm{C}$，需要的热量约是
\[200 \times 4.2 \times 80 \approx 6.7\times 10^{4}\ \mathrm{J}.\]
而把同样这杯水从 $100\,^{\circ}\mathrm{C}$ 全部烧成水蒸气，需要约 $200 \times 2250 \approx 4.5\times 10^{5}\ \mathrm{J}$——是前面那笔的七倍。烧开水时你早就见过这个事实：火力不变，水从室温到翻滚只要几分钟，可翻滚之后要很久很久才能烧干。大部分能量花在“拆网”（把分子从氢键网络里放走）上，而不是“加速”（升温）上。海边城市冬暖夏凉也是同一笔账：白天和夏天，海水吞下大量热量而温度只升一点；夜里和冬天，海水慢慢吐出热量。海洋就是沿海气候的巨大保温瓶。
\end{qiaoliang}

\subsection{罪状四：挑食}

盐放进水里会消失，油放进去却抱团分层。这两条其实都是那张网的“交友原则”。

先说盐。食盐是钠离子和氯离子的晶体（第 18 章）。水分子的负极端围住钠离子、正极端围住氯离子，一群水分子七手八脚把一个离子从晶格上“抬走”，裹成一个小水球——这叫\textbf{水合}。离子和水之间的吸引足以补偿拆散晶格、重排氢键网付出的代价，于是盐溶解了。糖虽然不带电荷，但分子里挂满亲水的羟基，能直接编进氢键网里，同样受欢迎。规律可以概括成四个字：\textbf{相似相溶}——极性的水，乐于接纳带电的或极性的伙伴。

这条规律在生活里极其好用。手上的机油、油漆用水洗不掉，用汽油或酒精一擦就掉——因为油污是非极性的碳氢链，只有同样“不显电性”的溶剂才接得住它；干洗店用的“干洗剂”其实就是不用水、改用非极性溶剂洗整件衣服。做饭也一样：盐、糖、醋往汤里一倒就化，因为它们亲水；而辣椒油里的辣味分子偏油溶，所以辣味要借油才能充分“跑”出来。

但“相似相溶”只是一句经验口诀，不是定律，它的局限同样值得知道。拿酒精家族来说：乙醇（酒精）和水能以任意比例互溶，因为它一端是亲水的羟基、碳链又很短；可是换成碳链更长的丁醇，就只能在水里溶解一部分——碳链越长，分子身上“油”的比重越大，水对它就越来越冷淡。到了六个碳以上的醇，就几乎和油一样被拒之门外。可见“亲不亲水”不是贴标签，而是分子内部两股势力——亲水部位和疏水部位——的比例之争。溶解的具体规则（饱和、温度、压强的影响）到第 25 章再细拆。

再看油。食用油分子的主体是长长的碳氢链，碳和氢的电负性几乎相同，整条链上没有明显的正负端，是个“电中性老好人”。把它放进水里会发生什么？注意，油分子和水分子之间并没有谁在“推开”谁——油和水的分子之间照样有普通的吸引（色散力），只不过比较弱。真正的障碍在水这一边：一个油分子混进氢键网，等于在网中间塞进一块牵不了手的异物，周围的水分子只好在油分子表面排成一圈拘谨的“笼子”，氢键的排列方式大大减少。让油分散开，等于强迫无数水分子放弃自由重排、集体列队——从统计上看，这种高度受限的状态极难自发出现（第 28 章会用“熵”把这个直觉说透）。于是能量和概率联手投票的结果，是水分子宁可彼此抱团、把油挤出去：油滴合并、分层，让受困的水分子最少。

所以请记住这句话：\textbf{疏水效应不是油分子之间存在什么特殊的“疏水键”}。油分子之间只有平平无奇的色散力；真正的驱动力是水的氢键网在“自保”。油不是被什么东西拉到一起的，是被水的社交网络“请出去”的。

\begin{shiyan}
\textbf{材料}：两个相同的小杯子、等量的水和食用油（各约半杯）、两支温度计（没有就用手指）、一个有阳光的窗台。

\textbf{步骤}：把两个杯子并排放在阳光下晒 30 分钟，每隔 10 分钟测一次（或指尖轻触杯壁比较）温度。回来后同样放置，比较谁先凉下来。

\textbf{观察点}：哪一杯升温快？哪一杯降温快？差别有多大？

\textbf{安全提示}：全程常温操作，无需加热；不要用玻璃温度计去测热水或开水；做完把油倒进厨余垃圾，不要倒进下水道。

你会看到食用油升温、降温都比水快。你刚才亲手摸到的，就是氢键网的“热惯性”——沿海城市的温柔气候和你的身体维持 $37\,^{\circ}\mathrm{C}$ 的能力，都是它的放大版。
\end{shiyan}

\section{现在，才轮到符号登场}

聊了这么久才写化学式，是故意的——先有意义，后有记号。

水分子的化学式是 \chem{H_2O}：一个氧，两个氢。要表达它的极性，化学家在两端正上方标上小小的希腊字母：
\[\chem{H^{\delta+}-O^{\delta-}-H^{\delta+}}\]
$\delta$（读作“德尔塔”）在这里表示“一点点”：$\delta+$ 是“带一丁点儿正电”，$\delta-$ 是“带一丁点儿负电”。注意它们\textbf{不是}完整的离子电荷——电子没有被抢走，只是在共享时偏向了氧。这和第 18 章里 \chem{Na^+}、\chem{Cl^-} 那种整个电子易主的离子，是两回事。

氢键则用三个点表示，画在两个分子之间：
\[\chem{O-H\cdots O}\]
实线是分子内部的共价键，三个点是分子之间的氢键。画水的结构图时，把这张网表示成“每个氧牵住四个邻居”的四面体格子就够了——但要记住，这只是人为的表示法：真实的氢键一刻不停地断开重连，而且分子本身没有任何颜色和“小棍子”。

\section{科学家怎么知道冰里面是空的}

\begin{zhengju}
冰浮在水上，这个现象人人都见过；可“为什么”曾经吵了很多年。最省事的猜测是：冰里面有气泡？可纯净的冰照样浮。也有人干脆说“固体就是会轻一些”——可几乎所有别的物质都不这样。要真正回答，必须\textbf{看见}冰里原子的排列。可原子太小了，任何显微镜都看不见，怎么办？

1912 年，德国物理学家劳厄想到一个办法：晶体的原子排列得整整齐齐，间距和 X 射线的波长差不多，晶体对 X 射线来说就像光栅对光一样，会产生\textbf{衍射}——穿过晶体后，X 射线在底片上打出一幅有规律的花样。反推这幅花样，就能算出原子的位置。布拉格父子随即把这套方法变成了实用的“原子照相机”。

把这台相机对准冰，谜底揭晓：冰中的氧原子排成一个个四面体搭成的网，网中间留着大空洞——冰确实“空旷”。还有一个更刁钻的问题：氢原子太轻，X 射线几乎“看”不见它，氢到底蹲在氧与氧之间的什么位置？后来问世的中子衍射对氢敏感，这才补上最后一块拼图：氢位于两个氧的连线上，而且偏向其中一边——正是我们今天画成 \chem{O-H\cdots O} 的那副样子。鲍林在 20 世纪 30 年代把这些线索整理成“氢键”概念，还顺手用它解释了一个更怪的观察：普通冰里的氢位置其实有点“乱”，这份乱留到绝对零度都冻结不掉，可以用氢在 \chem{O\cdots O} 之间两种摆法的可能数目定量算出来，和实验测得几乎分毫不差。

故事还没完。液态水的结构比冰难看得多——它一闪一闪地重排，抓不住定格。直到今天，“液态水里的氢键网到底是什么模样”仍然是个活跃的研究课题，不同实验给出的图像还在互相较劲。这正是科学的真实样子：冰浮不浮早已板上钉钉，而网的细节还在争论中前进。
\end{zhengju}

\section{这张网也有边界}

“极性分子＋氢键网”这个模型很好用，但别把它当成水的全部真相。

第一，\textbf{冰不止一种}。我们熟悉的冰是在常压下形成的；在不同的温度和压强下，水会结成结构完全不同的冰——科学家已经确认的冰晶型有十几种，有的比液态水还密，扔进水里会沉底。“冰一定浮”只是常压冰的性质，不是水的天性。

第二，\textbf{氢键是个有模糊地带的分类}。它的强弱跨越一个连续的范围，介于化学键和普通分子间作用之间，不同教材画的界线不完全一样。把它当成“强弱光谱上的一段”，比当成非黑即白的类别更接近事实。

第三，\textbf{别被商业词汇带跑}。市面上偶尔能看到“小分子团水”“活化水”“水的记忆”之类的宣传，声称经过某种处理的水更健康。氢键网确实在变化，但它的重排快得超乎想象——任何“结构”在万亿分之一秒的尺度上就已面目全非，不可能像刻录光盘一样保存信息。说水有“记忆”，就像说篝火能记住每一根木柴的形状。第 25 章还会看到，你喝下的水一进入身体，就立刻汇入全身的水，之前的一切“处理”都不会跟着走。

\begin{wuqu}
“油和水不相溶，是因为油分子之间有一种‘疏水键’把它们紧紧粘在一起。”——这个说法把因果关系彻底弄反了。反例就在眼前：把油滴分散在空气里（比如喷成油雾），油滴照样会聚拢成大滴——空气里可没有任何“疏水”的东西，油分子聚拢只是因为表面的分子被往内拉（第 22 章的表面张力）。而在水里，油分子之间也没有什么特殊胶水：驱使油抱团的，是水分子急于恢复自己的氢键网，把碍事的油挤出去。判断一种解释靠不靠谱，可以问一句：它把“推力”安在了谁身上？安在油身上，多半是误解；安在水的网络上，才是现代科学的答案。
\end{wuqu}

\section{回到那杯冰水}

开场看到的三件怪事，现在可以逐条结案了。

冰块浮在水上：冰的四面体脚手架留了空洞，密度只有液态水的约九成，十分之一露出水面——冰山一角，名副其实。

杯壁的水珠：空气里本来就漂着无数水分子。撞上冰冷的杯壁后，它们被“降速”，跑不出氢键网的抓捕，越聚越多，凝成你看见的水珠。水不是从杯子里漏出来的，是从空气里“收编”来的。

冰化水位不变：漂浮的冰块排开的水，恰好等于它融化后变成的水——冰块在水下占多大的“坑”，化掉就填多大的坑，所以水位几乎纹丝不动。（海冰融化不抬升海平面也是这个道理；真正让海平面上升的是陆地上的冰盖滑进海里。）

再加上那两个生活问题：温水擦身退烧，靠的是水蒸发时每克带走约 \ud{2400}{J} 的大手笔；海边冬暖夏凉，靠的是海水巨大的“热惯性”。一杯普通的冰水，从头到尾都是那张氢键网在表演。

\section{网的后面还有什么}

这一章你看到的是：\textbf{同一种作用力，在不同的组织方式下，给出了液体、固体、气候、体温这些尺度完全不同的现象}。水不是因为有哪一项属性特别强，而是因为四面体的氢键网“刚刚好”——强到能把分子留在液体里，又弱到一刻不停地重排，让水既能溶解盐糖，又能把油拒之门外。

故事有两个方向还没讲完。一个方向往下游：把盐、糖、气体乃至更多东西放进这张网里，网会怎样改造它们？溶解之后粒子在溶液里怎样排队、怎样导电、怎样隔着半透膜“较劲”？那是第 25 章《溶液里发生了什么》的领地。另一个方向往上游：如果一种分子一头亲水、一头亲油，把它扔进水里，水的“交友原则”会逼它摆出什么阵势？答案令人震惊——它们会自发围成空心球，把亲油的一头藏进球壁里。几十亿年前，正是这样一层自组装的水中“泡泡”，第一次把“里面”和“外面”分开，成了细胞的边界。第 51 章讲细胞膜时，我们会再次遇见今天这张网——到那时你会发现，生命的边界，本质上也是水“挑食”挑出来的。

\begin{tiaozhan}
水还有一个没讲的怪脾气：液态水的密度在 $4\,^{\circ}\mathrm{C}$ 时最大，无论升温还是降温都会变小。为什么？这是两股力量的拉锯。降温时，分子的热运动减弱，氢键更容易把邻居拉到“标准距离”排好，排列变整齐、间距变小，密度上升；但与此同时，越来越多的区域开始形成类似冰的四面体小团簇，团簇里的空洞又让体积膨胀。$4\,^{\circ}\mathrm{C}$ 附近，前一股力量刚好被后一股追平，密度到达顶点。再冷下去，空洞占了上风，水越接近结冰越“松”——所以 $0\,^{\circ}\mathrm{C}$ 的水已经比 $4\,^{\circ}\mathrm{C}$ 时轻，冬天湖底的水温恰恰维持在 $4\,^{\circ}\mathrm{C}$ 上下，成为水生生物的避难所。想再定量化一点，可以查这两个数：$4\,^{\circ}\mathrm{C}$ 时水的密度约 \ud{1.000}{g/cm^3}，$0\,^{\circ}\mathrm{C}$ 时约 \ud{0.9998}{g/cm^3}——差别很小，却足以驱动湖泊在入冬时整体对流一次，把氧气送到湖底。小数字，大生态。
\end{tiaozhan}

\section*{练一练}

\begin{sikaoti}
\item 画一个水分子，标出 $\delta+$ 和 $\delta-$ 端；再画上三个邻居，用实线表示共价键、用三个点表示氢键，让中间的分子恰好牵住四只“手”。图旁注明：这是帮助思考的表示法，真实的氢键在不断断开重连。
\item 氨（\chem{NH_3}）的分子也能形成氢键，但每个分子平均能牵住的邻居比水少。据此预测：氨的沸点和水相比如何？和甲烷相比呢？查数据验证。
\item 冬天，一条河流的表面结冰而河底不结冰。用本章的结构解释，画一幅“河流剖面图”，标出冰层、$4\,^{\circ}\mathrm{C}$ 水层和鱼的位置，并说明如果冰比水重会怎样。
\item 用“氢键网自保”的逻辑，一步步解释：为什么把油和水摇匀后，静置几分钟又会分层？要求不使用“油讨厌水”或“疏水键”这类说法，改从水分子一侧的能量与排列方式说起。
\item 1 克汗水蒸发约带走 \ud{2400}{J}，而 1 克水降低 $1\,^{\circ}\mathrm{C}$ 只放出约 \ud{4.2}{J}。估算：蒸发 10 克汗，理论上能把一个 60 千克的人（姑且把人体全当水）的体温拉低多少？再说说为什么实际效果比这个小。
\item 有人说：“某品牌处理过的水分子团更小，细胞更好吸收。”用本章关于氢键网重排速度的知识，写三句话反驳他。
\end{sikaoti}
